% OpenDC-Infer-LC — benchmark protocol / submission details (APPENDIX).
% Place after \appendix. Holds the spec machinery moved out of the main-body
% Methodology (Section~\ref{sec:methodology}) to fit the page budget: full corpus
% mechanics, the endpoint measurement contract, the cache-state taxonomy, the
% reference-harness procedure, and leaderboard/submission rules. Requires:
% \usepackage{enumitem}. Labels are unchanged so main-body cross-references resolve.

\section{Benchmark Protocol and Submission Details}
\label{sec:protocol_appendix}
This appendix gives the full specification deferred from the main-body methodology:
the synthetic-corpus generator, the endpoint measurement contract, the LC-Cache
state taxonomy, the reference-harness run procedure, and the leaderboard/submission
rules.

\subsection{Corpus Generation}
\label{subsec:corpus_details}
The long-context haystack is \textbf{procedurally generated from a seeded random
number generator}, never sourced from public documents. Each prompt follows a
RULER-style needle-in-a-haystack template: an instruction header, a synthetic
haystack built from a seeded sentence generator over fixed adjective/noun/verb/
connective vocabularies, one or more embedded \emph{needles}, and a question
footer. A needle is a distinctive ``magic number'' statement (\emph{``One of the
special magic numbers for \texttt{<key>} is: \texttt{<value>}.''}) with
\texttt{key} drawn from a fixed keyword set and \texttt{value} a random 7-digit
integer, giving an unambiguous exact-match target for the quality guardrail. We use
\textbf{single-key} (one needle) and \textbf{multi-key} (one queried needle plus
decoy needles with different keys at evenly spread depths) variants. A
length-control step re-tokenizes the decoded prompt with the declared tokenizer and
truncates/pads the filler so the total hits the target token count \emph{exactly}
($\leq8$-token drift, recorded per sample). Because token boundaries are
tokenizer-specific, the frozen sets are generated \emph{per tokenizer} (one tree
per model) and pinned by a \texttt{dataset\_version\_hash}. Synthetic generation is
chosen for three reasons: metric integrity (a repeated real passage is highly
compressible and would inflate prefix-cache hit rates, distorting the cache metrics
the benchmark measures), license cleanliness, and bit-for-bit reproducibility.

\subsection{Secondary Metrics (Full Set)}
\label{subsec:secondary_full}
Beyond SLO-qualified goodput (Eq.~\ref{eq:goodput}) and scaling efficiency
(Eq.~\ref{eq:efficiency}), the harness reports the following; Table~\ref{tab:metrics}
lists each with its availability marker.
\begin{itemize}[leftmargin=*, nosep, topsep=0pt, partopsep=0pt]
    \setlength{\itemsep}{0pt}\setlength{\parskip}{0pt}\setlength{\parsep}{0pt}\setlength{\topsep}{0pt}
    \item \textbf{Latency:} TTFT, TPOT, and E2E at p50/p95/p99.
    \item \textbf{Throughput:} raw output tok/s, total tok/s, and requests/s.
    \item \textbf{Capacity:} max concurrent users and max request rate under the SLO.
    \item \textbf{Reliability:} success, timeout, server-error, cancellation, and OOM rates.
    \item \textbf{Cache behavior:} prefix-cache hit rate, KV-cache usage, placement, eviction policy, and cache-reuse speedup.
    \item \textbf{Transport:} declared/observed transport, fallback count, and cross-node communication evidence.
    \item \textbf{Storage:} model/prompt/cache/log paths, read/write bandwidth, and whether each path is in the timed request path.
    \item \textbf{Hardware efficiency:} output tok/s per accelerator and, when available, output tokens per joule.
    \item \textbf{Run stability:} run-to-run variance, p99 drift, thermal-throttling indicators, telemetry gaps, and post-saturation behavior.
\end{itemize}

\subsection{Endpoint Measurement Contract}
\label{subsec:endpoint_contract}
The reference harness drives an OpenAI-compatible streaming endpoint (or
equivalent). Because endpoints differ in streaming behavior, we fix:
\begin{itemize}
    \item \textbf{Request start:} immediately before the client sends the request.
    \item \textbf{TTFT:} start until the first \emph{non-empty} generated token; metadata-only or empty chunks do not count.
    \item \textbf{TPOT:} $\mathrm{TPOT}_i = (t^{\mathrm{end}}_i - t^{\mathrm{first}}_i)/\max(y_i-1,1)$, with $y_i$ output tokens and $t^{\mathrm{first}}_i,t^{\mathrm{end}}_i$ the first-token and completion timestamps.
    \item \textbf{E2E:} start until the final completion event is parsed and the stream closes cleanly.
    \item \textbf{Token accounting:} prompts/outputs re-tokenized with the declared tokenizer unless a validated server-side contract is provided.
    \item \textbf{Timeouts/cancellations:} timed-out, cancelled, malformed, or incomplete requests are unsuccessful and count against the success rate.
    \item \textbf{Retries:} disabled by default; if enabled for a diagnostic run, reported separately.
\end{itemize}
The submission package includes a \texttt{loadgen\_config.yaml} specifying
load-generator hosts, CPU count, NICs, network path, client processes, connection
settings, TLS mode, tokenizer mode, timeouts, retry policy, and clock source. The
load generator must not be the bottleneck at the reported peak: submissions report
host count, CPU/NIC utilization, client process count, connection/timeout settings,
tokenizer mode, and clock source, and a run in which the generator saturates CPU,
network, file I/O, TLS, connection pools, or tokenization is invalid for ranking.

\subsection{LC-Cache State Rules}
\label{subsec:cache_rules}
Because cache state strongly affects results, cache modes are reported separately:
\begin{itemize}
    \item \textbf{Cold cache:} no prefix/KV/semantic/document/storage cache populated before the timed window.
    \item \textbf{Benchmark-warmed cache:} reuse arises only from requests issued during the timed window.
    \item \textbf{Prefilled cache:} the shared prefix / reusable state is loaded before the timed window.
    \item \textbf{Persistent session cache:} state persists across turns within one synthetic conversation.
    \item \textbf{Cross-run cache:} state survives across independent runs---\emph{prohibited} for ranked submissions unless disclosed in a separate non-default category.
\end{itemize}
Each LC-Cache submission discloses whether prefix caching, paged KV cache, KV
offload, semantic cache, document cache, or storage-backed cache is enabled, and
reports cache key granularity, capacity, eviction policy, observed hit rate, reset
procedure, and memory location (HBM, host DRAM, local NVMe, shared filesystem, or
remote object storage). LC-Cache leaderboards are reported separately for
cold-, benchmark-warmed-, and prefilled-cache modes.

\subsection{Reference Harness Procedure}
\label{subsec:harness}
The reference harness handles workload generation, request scheduling,
streaming-response measurement, token accounting, SLO validation, quality
checking, metadata collection, and report generation. It is backend-neutral and
talks to any endpoint satisfying the measurement contract. A run proceeds:
\begin{enumerate}
    \item Select workload, model tier, SLO profile, cache mode, and load-generation mode.
    \item Verify the system passed the data-center qualification phase (Section~\ref{sec:execution_protocol}).
    \item Warm up the serving system under a fixed warmup policy (including a priming request that absorbs first-request kernel/graph compilation).
    \item Run the timed measurement window at a fixed offered load.
    \item Record per-request metrics, streaming timestamps, system metrics, transport logs, and load-generator metrics.
    \item Compute latency, reliability, raw throughput, SLO-qualification status, and SLO-qualified goodput.
    \item Run the quality guardrail on the validation subset.
    \item Emit the standardized submission package and report.
\end{enumerate}

\subsection{Leaderboard and Submission}
\label{subsec:leaderboard}
OpenDC-Infer-LC reports multiple leaderboards rather than one aggregate score. The
main leaderboard ranks by SLO-qualified goodput on LC-32K (the primary enterprise
RAG setting); additional leaderboards cover LC-8K interactive capacity, LC-128K
heavy long-context capacity, LC-Cache prefix-reuse efficiency, and hardware
efficiency. A submission is ranked only if: the run is SLO-qualified; the quality
guardrail is satisfied; the system qualification phase passes; required metadata
and raw artifacts are provided; and the cache-state, endpoint-measurement, and
timed-window rules are followed. Submissions lacking complete qualification
evidence may be listed as preliminary or partially qualified but are not ranked
against fully qualified ones.
